\newif\ifglobal

%%%%%%%%%%%%%%%%%%%%%%%
%%% Compile to view %%%
%%%%%%%%%%%%%%%%%%%%%%%

%%%%%%%%%%%%%%%%%%%%%%%%%%%%%%%%%%%%%%%%%%%%%%%%%%%
%%% Readme:                                     %%%
%%% https://www.overleaf.com/read/bwdjvfjksvjy/ %%%
%%%%%%%%%%%%%%%%%%%%%%%%%%%%%%%%%%%%%%%%%%%%%%%%%%%

% >>> M?: marks a module setting number or important notice

% >>> M4: Set {./relative-path} to external files for this module <<<
\def \modulefiles {./41-PARLIO}
% To add an external file, use:
% (Replace '---' with actual filename)
% '\input{\modulefiles/tables/---.tex}' for tables
% '\includegraphics{\modulefiles/figures/---}' for figures
% '\subfile{\modulefiles/---}' for register files


% >>> M1: This module is in global project? <<<
% 'YES' - uncomment; 'NO' - comment out
\globaltrue

\ifglobal
    \documentclass[main\_\_CN.tex]{subfiles}
   
\else
    % Font "FandolSong-Regular" used by ctex does not have GSUB table.
    % It causes warnings that do not affect compilation and can be ignored.
    \PassOptionsToPackage{quiet}{fontspec}

    \documentclass[openany, 10pt]{book}
   
    % >>> M2: In ./00-shared/preamble-trm-module, also find
    %         settings for visibility of todo notes, watermark <<< 
    \usepackage{preamble-trm-module}


    % Enable Chinese typesetting
    \usepackage{ctex}

    % >>> M3: Temporary labels in Chinese
    %         you can add your own labels there <<<
    \myexternaldocument{./00-shared/temp-labels-cn}

\fi %\ifglobal


\setcounter{secnumdepth}{4}


% >>> M5: If this module needs any updates in
% './00-shared/preamble-trm-module.sty'
% (include additional package, etc.), contact your module PM
% or create the file './00-shared/changelog.tex' make the
% necessary updates yourself and document them in this file <<<


\begin{document}


% Sets language into which pre-defined bits of text are translated
\selectlanguage{chinese}

% Do not delete this block
\ifglobal
% Add the following front matter if
% this module is outside of global project
\else
    \InputIfFileExists{readme}{}{}
    {\let\clearpage\relax \listoftodos}
    {\let\clearpage\relax \listofdonetodos}
    \tableofcontents
    %\listoftables
    %\listoffigures
\fi


%%%%%%%%%%%%%%%%%%%%%%%%%
%%% TRM Module Begins %%%
%%%%%%%%%%%%%%%%%%%%%%%%%

\hypertarget{parlio}{}
\chapter{并行 IO 控制器 (PARL\_IO)}
\label{mod:parlio}

%%%%%%%%%%%%%%%%%%%%
%%% Introduction %%%
%%%%%%%%%%%%%%%%%%%%

\section{概况}

\chipname{} 包含一个并行 IO 控制器 (PARLIO)，支持通过 GDMA 在并行总线上 实现外部设备和内部存储器之间的数据通信。它由 TX 和 RX 两个子模块组成，分别作为发送器和接收器。PARLIO 支持全双工通信。

鉴于 IO 数据的灵活性，PARLIO 可用作通用接口，连接各种外围设备。例如可以把 SPI 作为主机，PARLIO 作为从机，实现点对点传输。详细的应用实例，请参考小节 \ref{sec:parlio-application-example}。

\section{术语}
为了更好地说明 PARLIO 的功能，本章使用了以下术语。

\begin{longtable}[c]{ R{4cm} p{10cm} }

\textbf{RX 模块}      & PARLIO 子模块，从外部并行总线接收数据并将其存储在内部存储器中。 \\
\textbf{TX 模块}      & PARLIO 子模块，将数据从内部存储器传输到外部并行总线上。 \\
\textbf{RXD}          & 从 RX 模块的 IO 接口接收的并行数据。 \\
\textbf{TXD}          & 从 TX 模块 IO 接口发送的并行数据。  \\
\textbf{帧}           & 传输数据的单位，从设置 START 信号时起，到收到 EOF 信号止。 \\
\textbf{自由运行时钟} & 持续翻转的时钟。若非自由运行时钟，则时钟只在收到有效数据的期间翻转，其余时间停止翻转。\\
\textbf{GDMA SUC EOF} & GDMA 帧传输成功结束的信号。GDMA 接收到此信号时，触发一个 GDMA 中断，表明当前帧正确，且数据接收结束。 \\
\textbf{GDMA ERR EOF} & GDMA 帧传输错误结束的信号。GDMA 接收到此信号时，触发一个 GDMA 中断，表明当前帧错误，且数据接收结束。 \\
\textbf{CDC}          & 跨时钟域处理 (Clock Domain Crossing)。 \\

\end{longtable}

%%%%%%%%%%%%%%%%%%%%
%%% Feature List %%%
%%%%%%%%%%%%%%%%%%%%

\section{特性}
\label{sec:parlio-features}

PARLIO 模块具备以下特性：

\begin{itemize}
    \item 支持多种时钟源可选：
    \begin{itemize}
        \item 包括外部 IO 时钟 PAD\_CLK、内部系统时钟 XTAL\_CLK、PLL\_F240M\_CLK 和 RC\_FAST\_CLK
        \item 支持最大时钟频率 40 MHz
        \item 支持对时钟频率进行整数分频
    \end{itemize}
    \item 支持传输数据总线位宽配置为 1/2/4/8/16 位
    \item 支持 16 位半双工传输和 8 位全双工传输
    \item 在 1/2/4 位模式下，支持在一个字节范围内对比特数据重新排序
    \item 包含用于接收 IO 并行数据的 RX 子模块：
    \begin{itemize}
        \item 支持 RX 子模块输入时钟取反
        \item 支持多种接收模式
        \item 支持 GDMA SUC EOF 信号生成模式可配置
        \item 支持外部使能信号的 IO 管脚可配置
    \end{itemize}
    \item 包含用于发送 IO 并行数据的 TX 子模块：
    \begin{itemize}
        \item 支持 TX 子模块输出时钟取反
        \item 支持有效信号输出
        \item 支持总线空闲时数值可配置
    \end{itemize}

\end{itemize}


%%%%%%%%%%%%%%%%%%%%%%%%%%%%%%
%%% Architectural Overview %%%
%%%%%%%%%%%%%%%%%%%%%%%%%%%%%%

\section{系统架构}
\label{sec:parlio-arch-overview}

\begin{figure}[H]
    \centering
    \includegraphics[width=1\textwidth]{\modulefiles/figures/41-parlio_architecture_20230217.pdf}
    \caption{PARLIO 系统框图}
    \label{fig:parlio-architecture}
\end{figure}

图 \ref{fig:parlio-architecture} 展示了 PARLIO 的系统内部架构。除了 RX 模块和 TX 模块外，PARLIO 还包括一组状态配置寄存器。

RX 模块将 RXD 转换为异步 FIFO 接口，异步 FIFO 将 RXD 同步到 AHB 时钟域，最后 RXD 被转换为标准的 GDMA 接口，并被传输到内部存储器。

TX 模块通过 GDMA 从内部存储器获取数据，并将 GDMA 接口转换为异步 FIFO 接口。异步 FIFO 将数据同步到 TX Core 时钟域，最后将其转换为 TXD，用于并行 IO 总线输出。

%%%%%%%%%%%%%%%%%%%%%%%%%%%%%%
%%% Functional Description %%%
%%%%%%%%%%%%%%%%%%%%%%%%%%%%%%

\section{功能描述}
\label{sec:parlio-func-descr}

\subsection{时钟生成器}
\label{subsec:PARLIO-clock_cfg}

PARLIO 模块有四个输入时钟域，即 RX Core、TX Core、AHB 和 APB。

状态配置寄存器组位于 APB 时钟域内。

GDMA 接口逻辑位于 AHB 时钟域内。

RX Core 和 TX Core 时钟域各有四个时钟源可供选择，分别为内部系统时钟源 XTAL\_CLK、RC\_FAST\_CLK、PLL\_F240M\_CLK 和外部时钟源 (PAD\_CLK\_TX/RX)，如图 \ref{fig:parlio-clock-gen} 所示。配置 \hyperref[fielddesc:PCRPARLCLKRXSEL]{PCR\_PARL\_CLK\_RX\_SEL} 和 \hyperref[fielddesc:PCRPARLCLKTXSEL]{PCR\_PARL\_CLK\_TX\_SEL} 选择时钟源，配置 \hyperref[fielddesc:PCRPARLCLKRXDIVNUM]{PCR\_PARL\_CLK\_RX\_DIV\_NUM} 和 \hyperref[fielddesc:PCRPARLCLKTXDIVNUM]{PCR\_PARL\_CLK\_TX\_DIV\_NUM} 对时钟进行分频，分频系数最大可配置为 \begin{math} (2^{16}-1) \end{math}。

RX 模块支持对输入时钟进行取反。TX 模块的工作时钟可以输出到 IO，并在输出前支持对输出时钟进行取反。

\begin{figure}[H]
    \centering
    \includegraphics[width=0.7\textwidth]{\modulefiles/figures/41-parlio-clock-gen_20230106.pdf}
    \caption{PARLIO 时钟生成}
    \label{fig:parlio-clock-gen}
\end{figure}


\subsection{时钟复位使用限制}
\label{subsec:parlio-clock_reset_restriction}

由于 PARLIO 模块的通用性，其 PAD 时钟可能来自不同的主机（外部设备或内部时钟源），这些时钟可能是自由运行时钟，也可能不是。如果外部运行时钟不是自由运行时钟，PARLIO 内部的一些控制信号将无法进行跨时钟域处理，因此对操作过程会有一定限制。

\begin{enumerate}
    \item 异步 FIFO 复位时，需要两个时钟周期才能从 AHB 时钟域同步到 Core 时钟域。因此如果需要在使用非自由运行的时钟时复位 AHB 时钟域，需保证提前两个时钟周期进行复位同步。具体的操作步骤为：
    \begin{itemize}
        \item \textbf{场景一：}当前帧传输基于自由运行时钟，下一帧传输不是基于自由运行时钟。\\
        \textbf{具体操作：}用户可以在切换到非自由运行时钟之前复位下一帧传输，并在复位完成后配置时钟切换。
        \item \textbf{场景二：}当前帧传输不是基于自由运行时钟，下一帧传输基于自由运行时钟。\\
        \textbf{具体操作：}此时只需确保在复位和正式开始传输期间有两个时钟周期的间隔。
        \item \textbf{场景三：}当前帧传输不是基于自由运行时钟，下一帧传输也不是基于自由运行时钟。\\
        \textbf{具体操作：}如需复位下一帧传输，用户需要先切换到内部的自由运行时钟，复位完成后再切换回实际时钟。
    \end{itemize}
    
    \item 由于非自由运行时钟造成的限制，\hyperref[fielddesc:PARLIORXSTART]{PARL\_IO\_RX\_START} 和 \hyperref[fielddesc:PARLIOTXSTART]{PARL\_IO\_TX\_START} 无法进行跨时钟域处理，所以在操作时必须等到 \hyperref[fielddesc:PARLIORXSTART]{PARL\_IO\_RX\_START} 和 \hyperref[fielddesc:PARLIOTXSTART]{PARL\_IO\_TX\_START} 稳定后再开始数据传输，否则可能会发生亚稳态的问题。
    
    RX 模块中的具体操作步骤为：
    \begin{itemize}
        \item 清除 \hyperref[fielddesc:PCRPARLCLKRXEN]{PCR\_PARL\_CLK\_RX\_EN} 关闭 RX Core 时钟域；
        \item 向 \hyperref[fielddesc:PARLIORXSTART]{PARL\_IO\_RX\_START} 写 1；
        \item 置位 \hyperref[fielddesc:PCRPARLCLKRXEN]{PCR\_PARL\_CLK\_RX\_EN} 打开 RX Core 时钟域；
        \item 操作外部设备，使其开始发送数据；
        \item 清除 \hyperref[fielddesc:PCRPARLCLKRXEN]{PCR\_PARL\_CLK\_RX\_EN} 关闭 RX Core 时钟域；
        \item 向 \hyperref[fielddesc:PARLIOTXSTART]{PARL\_IO\_TX\_START} 写 0。
    \end{itemize}
    TX 模块中的具体操作步骤为：
    \begin{itemize}
        \item 清除 \hyperref[fielddesc:PCRPARLCLKTXEN]{PCR\_PARL\_CLK\_TX\_EN} 关闭 TX Core 时钟域；
        \item 向 \hyperref[fielddesc:PARLIOTXSTART]{PARL\_IO\_TX\_START} 写 1；
        \item 置位 \hyperref[fielddesc:PCRPARLCLKTXEN]{PCR\_PARL\_CLK\_TX\_EN} 打开 TX Core 时钟域；
        \item 操作外部设备，使其开始接收数据；
        \item 清除 \hyperref[fielddesc:PCRPARLCLKTXEN]{PCR\_PARL\_CLK\_TX\_EN} 关闭 TX Core 时钟域；
        \item 向 \hyperref[fielddesc:PARLIOTXSTART]{PARL\_IO\_TX\_START} 写 0。
    \end{itemize}
    
    \item 复位需遵循以下要求：
    \begin{itemize}
        \item 芯片启动时，时钟复位应遵循以下顺序：
        \begin{itemize}
            \item 首先复位 APB 时钟域；
            \item 其次复位 AHB 时钟域；
            \item 最后复位 Core 时钟域。
        \end{itemize}
        \item 帧间传输需要复位 Core 时钟域和异步 FIFO。
    \end{itemize}
\end{enumerate}

\subsection{主从机模式}

TX 模块支持主从机模式，而 RX 模块只支持从机模式。

当 TX 模块作为主机时，必须将时钟源设置为内部自由运行时钟。TX 模块在时钟的上升沿驱动 TXD。

当 TX 模块作为从机时，有以下三种情况：

\begin{itemize}
    \item \textbf{场景一：}主机发送的时钟为自由运行时钟。\\
    \textbf{要求：}此时对主机的驱动时钟采集边沿没有要求。
    \item \textbf{场景二：}主机发送的时钟非自由运行时钟，且其时钟波形如图 \ref{fig:parlio-pos-clock} 所示。\\
    \textbf{要求：}此时主机需要在下降沿时采集。
    \begin{figure}[H]
    \centering
    \includegraphics[width=0.8\textwidth]{\modulefiles/figures/41-parlio-pos-clock.png}
    \caption{主机时钟正向波形}
    \label{fig:parlio-pos-clock}
    \end{figure}
    \item \textbf{场景三：}主机发送的时钟非自由运行时钟，且其时钟波形如图 \ref{fig:parlio-neg-clock} 所示。\\
    \textbf{要求：}此时主机在输出前需要对原时钟进行取反，将其转换成图 \ref{fig:parlio-pos-clock} 所示的波形。
    \begin{figure}[H]
    \centering
    
\includegraphics[width=0.8\textwidth]{\modulefiles/figures/41-parlio-neg-clock.png}
    \caption{主机时钟负向波形}
    \label{fig:parlio-neg-clock}
    \end{figure}
\end{itemize}

RX 模块只能作为从机，可能出现以下三种情况：
\begin{itemize}
    \item \textbf{场景一：}主机发送的时钟是自由运行时钟。\\
    \textbf{要求：}此时对主机的驱动时钟采集边沿没有要求，采集的有效数据以外部使能信号为准。
    \item \textbf{场景二：}主机发送的时钟非自由运行时钟，且其时钟波形如图 \ref{fig:parlio-pos-clock} 所示。\\
    \textbf{要求：}此时，主机需在上升沿驱动数据，而 RX 模块需在下降沿采集数据（即对主时钟进行翻转）。
    \item \textbf{场景三：}主机发送的时钟非自由运行时钟，且其时钟波形如图 \ref{fig:parlio-neg-clock} 所示。\\
    \textbf{要求：}此时，主机需在下降沿驱动数据，而 RX 模块需在上升沿采集数据（即使用原始主时钟）。
\end{itemize}

\subsection{RX 模块接收模式}

PARLIO 的 RX 模块有 15 种接收模式，可根据使能信号规格分为三大类：

\begin{itemize}
    \item 电平使能模式：接收的数据由外部信号电平使能；
    \item 脉冲使能模式：接收的数据由外部信号脉冲使能；
    \item 软件使能模式：数据接收的使能信号可由用户直接配置。
\end{itemize}


\subsubsection{电平使能模式}

根据外部使能信号的有效电平，电平使能模式可分为两种子模式，如图 \ref{fig:parlio-level-sub-mode} 所示。 

在这两种子模式下，外部使能信号的有效电平必须与有效数据对齐。由于外部电平使能信号需要占用一个 IO 管脚，RXD 可用的 IO 管脚最多为 15 个。

\begin{figure}[H]
\centering
\includegraphics[width=0.8\textwidth]{\modulefiles/figures/41-parlio-level-submode}
\caption{RX 模块电平使能子模式}
\label{fig:parlio-level-sub-mode}
\end{figure}

\subsubsection{脉冲使能模式}
根据脉冲有效电平及其与有效数据的对应性，脉冲使能模式可分为 12 种子模式，详细分类可参见图 \ref{fig:parlio-pulse-sub-mode}。

子模式 1 \verb+~+ 8 都包含起始脉冲和结束脉冲，不同之处在于起始脉冲和结束脉冲是否与有效数据对齐。

子模式 9 \verb+~+ 12 只包含起始脉冲，通过配置 \hyperref[fielddesc:PARLIORXDATABYTELEN]{PARL\_IO\_RX\_DATA\_BYTELEN} 可表示有效数据传输结束。

由于外部脉冲使能信号需要占用一个 IO 管脚，RXD 可用的 IO 管脚最多为 15 个。但在子模式 4、8、10、12 下，数据在脉冲的第一个边沿前和最后一个边沿后被认为有效，可以实现使能信号的 IO 管脚与数据 IO 管脚的复用，此时 RXD 可使用 16 个 IO 管脚。

\begin{figure}[H]
\centering
\includegraphics[width=0.8\textwidth]{\modulefiles/figures/41-parlio-pulse-submode}
\caption{RX 模块脉冲使能子模式}
\label{fig:parlio-pulse-sub-mode}
\end{figure}


\subsubsection{软件使能模式}

软件使能模式下的使能信号由内部配置寄存器决定。用户切换到此模式时，只有同时写 1 到 \hyperref[fielddesc:PARLIORXSWEN]{PARL\_IO\_RX\_SW\_EN} 和 \hyperref[fielddesc:PARLIORXSTART]{PARL\_IO\_RX\_START}，才能开始接收数据。

由于使能信号不占用接口 IO 管脚，RXD 在接口上的最大可用 IO 管脚数量为 16。由于时钟域不同，使能信号无法与有效数据对齐，因此只能通过有效时钟沿来判定有效数据。此时，RX Core 时钟需要与有效数据一一对应。

\begin{figure}[H]
\centering
\includegraphics[width=0.8\textwidth]{\modulefiles/figures/41-parlio-software-submode}
\caption{RX 模块软件使能子模式}
\label{fig:parlio-software-sub-mode}
\end{figure}

\subsection{RX 模块 GDMA SUC EOF 信号生成}

RX 模块产生一个 GDMA SUC EOF 信号，表示当前帧传输结束，并将其发送到 GDMA 接口。GDMA SUC EOF 信号可由外部使能信号或内部配置字节长度产生。

\begin{itemize}
    \item 当 GDMA SUC EOF 信号由内部配置的字节长度生成时，接收模式的选择不受限制。但在这种情况下必须配置 \hyperref[fielddesc:PARLIORXDATABYTELEN]{PARL\_IO\_RX\_DATA\_BYTELEN}。如果 \hyperref[fielddesc:PARLIORXDATABYTELEN]{PARL\_IO\_RX\_DATA\_BYTELEN} 的配置值小于实际接收的数据，将提前触发 GDMA SUC EOF。此时 RX 模块停止从 FIFO 读取数据，但 FIFO 会继续接收外部数据，直到 RX\_FIFO\_WFULL\_INTR 中断被触发。
    \item 当 GDMA SUC EOF 信号由外部使能信号产生时，只可选择脉冲使能模式的子模式 1、3、5 和 7。在这种模式下，传输不受 \hyperref[fielddesc:PARLIORXDATABYTELEN]{PARL\_IO\_RX\_DATA\_BYTELEN} 值的影响，并且传输帧的数据长度不受限制。
\end{itemize}

\subsection{RX 模块超时}

RX 模块支持接收超时功能。当超时被触发时，产生一个 GDMA ERR EOF 信号，并发送至 GDMA 接口以表示结束数据接收。配置 \hyperref[fielddesc:PARLIORXTIMEOUTTHRESHOLD] {PARL\_IO\_RX\_TIMEOUT\_THRESHOLD} 可设置超时阈值。

默认情况下启用超时功能，用户也可以选择禁用该功能。超时阈值的配置上限为 AHB 时钟域的 \begin{math} (2^{16}-1) \end{math} 周期，下限取决于 AHB 时钟域和 RX Core 时钟域之间的相对频率关系。建议将 \hyperref[fielddesc:PARLIORXTIMEOUTTHRESHOLD] {PARL\_IO\_RX\_TIMEOUT\_THRESHOLD} 设置为较大的数值，以避免生成非期望的 GDMA ERR EOF 信号。


\subsection{TX 模块有效信号输出}

TX 模块可生成与 TXD 对齐的有效信号，可通过配置 \hyperref[fielddesc:PARLIOTXHWVALIDEN]{PARL\_IO\_TX\_HW\_VALID\_EN} 选择是否将其输出到 TXD。有效信号的极性固定为高电平有效。

默认情况下禁用有效输出功能。若启用该功能，输出的有效信号将占据 TXD 的最高有效位，即无论原始值为多少，TXD 的第 15 位将维持在高电平，并作为有效信号输出。然而，有效信号管脚并不影响总线宽度的配置。例如，如果用户将数据总线宽度配置为 1 位，即数据管脚为 TXD[0]，此时用户仍然可以启用有效输出功能，将 TXD[15] 固定为有效信号输出管脚。

\subsection{TX 模块总线空闲值}
TX 模块不传输数据时被认为保持在空闲状态。TX 模块支持总线空闲值可配置。

总线空闲值默认为 0x0，最大可配置值为 0xFFFF。注意，应避免配置的空闲值与其他已开启的功能产生冲突，例如当使用 TXD 的最高有效位作为有效信号时，应避免将空闲值的最高有效位配置为 1。

\subsection{单帧数据传输}

RX/TX 模块以字节为单位进行传输，即单帧至少传输 1 个字节的数据。

当 RX 模块通过字节长度生成 GDMA EOF 信号时，单帧传输的最大长度是 \begin{math} (2^{16}-1) \end{math} 字节。当 RX 模块通过外部设备的使能信号产生 GDMA EOF 信号时，单帧传输的字节数没有限制。

TX 模块仅通过字节长度生成 GDMA EOF 信号，因此单帧传输的最大长度为 \begin{math} (2^{16}-1) \end{math} 字节。

当配置的数据总线宽度为 16 位时，字节长度必须配置为 2 字节的倍数。

通常情况下，PARLIO 可以进行全双工传输，但如果处于 16 位总线宽度模式，由于 IO 数量限制，PARAIO 只能进行半双工传输。

\subsection{字节范围内数据重排}

一个字节范围内的数据顺序可以重新排序。以 RX 模块为例，当配置总线位宽为 2 比特时，数据在被写入 RX FIFO 之前需要被打包成一个字节。

假设原始的比特数据排列顺序为：

\begin{center}
\{ \{ b\_{0}, b\_{1} \}, \{ b\_{2}, b\_{3} \}, \{ b\_{4}, b\_{5} \}, \{ b\_{6}, b\_{7} \} \}
\end{center}

则启用该功能后，数据顺序将更改为：

\begin{center}
\{ \{ b\_{6}, b\_{7} \}, \{ b\_{4}, b\_{5} \}, \{ b\_{2}, b\_{3} \}, \{ b\_{0}, b\_{1} \} \}
\end{center}

%%%%%%%%%%%%%%%%%%%%%%%%%%%%%%
%%% Programming Procedures %%%
%%%%%%%%%%%%%%%%%%%%%%%%%%%%%%

\section{配置流程}
\label{sec:parlio-prog-procedures}

\subsection{数据接收操作流程}

本节介绍 RX 模块接收数据的配置流程。参照以下流程，从连接到外部设备的 IO 管脚接收并行数据，并存储在内部存储器中。关于 RX 模块时钟和复位的详细操作限制，请参考小节 \ref{subsec:parlio-clock_reset_restriction}。

\begin{enumerate}
    \item 复位 RX 模块。具体复位方法与顺序，请参考小节 \ref{subsec:parlio-clock_reset_restriction}。
    \item 置位 \hyperref[fielddesc:PARLIORXFIFOWFULLINTCLR]{PARL\_IO\_RX\_FIFO\_WFULL\_INT\_CLR} 和 \hyperref[fielddesc:PARLIORXFIFOWFULLINTENA]{PARL\_IO\_RX\_FIFO\_WFULL\_INT\_ENA}。
    \item 选择 RXD IO 管脚。如果用户使用 PAD 时钟，则也需配置时钟 IO 管脚。
    \item 配置 PCR 寄存器，选择时钟源并对时钟进行分频。
    \item 关闭 RX 模块 Core 时钟域的时钟。
    \item 按照小节 \ref{sec:parlio-features} 和 \ref{sec:parlio-func-descr} 中的描述，选择所需接收模式并启用相关功能。
    \item 配置 GDMA 接收链表。
    \item 置位 \hyperref[fielddesc:PARLIORXREGUPDATE]{PARL\_IO\_RX\_REG\_UPDATE}，同步寄存器信号。
    \item 置位 \hyperref[fielddesc:PARLIORXSTART]{PARL\_IO\_RX\_START}。
    \item 打开 RX 模块 Core 时钟域的时钟。
    \item 操作外部设备，使其发送数据。
    \item 轮询 GDMA SUC EOF 中断信号。
    \item 清除 GDMA SUC EOF 中断信号。
    \item 关闭 RX 模块 Core 时钟域的时钟。
    \item 清除 \hyperref[fielddesc:PARLIORXSTART]{PARL\_IO\_RX\_START}。
\end{enumerate}

\subsection{数据发送操作流程}

本节介绍 TX 模块发送数据的编程流程。参照以下流程，将内部存储器中的并行数据传输到与外部设备相连的 IO 管脚。关于 TX 模块时钟和复位的详细操作限制，请参考小节 \ref{subsec:parlio-clock_reset_restriction}。

\begin{enumerate}
    \item 复位 TX 模块。具体复位方法与顺序，请参考小节 \ref{subsec:parlio-clock_reset_restriction}。
    \item 依次置位 \hyperref[fielddesc:PARLIOTXFIFOREMPTYINTCLR]{PARL\_IO\_TX\_FIFO\_REMPTY\_INT\_CLR}、\hyperref[fielddesc:PARLIOTXEOFINTCLR]{PARL\_IO\_TX\_EOF\_INT\_CLR}、\hyperref[fielddesc:PARLIOTXFIFOREMPTYINTENA]{PARL\_IO\_TX\_FIFO\_REMPTY\_INT\_ENA} 和 \hyperref[fielddesc:PARLIOTXEOFINTENA]{PARL\_IO\_TX\_EOF\_INT\_ENA}。
    \item 选择 TXD IO 管脚。如果用户使用 PAD 时钟，则也需配置时钟 IO 管脚。
    \item 配置 PCR 寄存器，选择时钟源并对时钟进行分频。
    \item 关闭 TX 模块 Core 时钟域的时钟。
    \item 按照小节 \ref{sec:parlio-func-descr} 中的描述，选择开启所需功能。
    \item 配置 GDMA 发送链表。
    \item 轮询 \hyperref[fielddesc:PARLIOTXREADY]{PARL\_IO\_TX\_READY}。
    \item 置位 \hyperref[fielddesc:PARLIOTXSTART]{PARL\_IO\_TX\_START}。
    \item 打开 TX 模块 Core 时钟域的时钟。
    \item 操作外部设备，使其接收数据。
    \item 轮询 \hyperref[fielddesc:PARLIOTXEOFINTST]{PARL\_IO\_TX\_EOF\_INT\_ST} 中断信号。
    \item 置位 \hyperref[fielddesc:PARLIOTXEOFINTCLR]{PARL\_IO\_TX\_EOF\_INT\_CLR} 中断信号。
    \item 关闭 TX 模块 Core 时钟域的时钟。
    \item 清除 \hyperref[fielddesc:PARLIOTXSTART]{PARL\_IO\_TX\_START}。
\end{enumerate}


\section{应用示例}
\label{sec:parlio-application-example}

本节介绍 PARLIO 模块的应用示例及其具体操作过程。示例中使用的外围设备都来自 ESP 系列芯片，这些外设均可以和 PARLIO 模块组成完整的数据通路。

\begin{tiplisting}
示例中所构建的数据通路可能不是最佳方案，例如，在实际应用中，用户可使用两颗相同 ESP 芯片上的 SPI 外设来完成点对点传输，而无需使用 PARLIO。但本小节中的示例在一定程度上证明了 PARLIO 接口的灵活性。\end{tiplisting}

\subsection{与 SPI 进行数据传输}

在本示例中，外部 SPI 作为主机发送数据，PARLIO RX 模块作为从机接收数据，与此同时，PARLIO TX 模块作为主机发送数据，外部 SPI 作为从机接收数据，从而实现点对点串行数据传输。

\begin{itemize}
    \item 实现 SPI 发送、PARLIO 接收，请参考以下操作过程：
    \begin{itemize}
        \item 配置 SPI 时钟。
        \item 将 SPI 配置为主机。
        \item 配置信号管脚，将 FSPICLK 管脚连接至 PAD\_CLK\_RX，FSPICS0 管脚连接至 RXD[16]，FSPID 管脚连接至 RXD[0]。 
        \item 将发送的数据写入 SPI 缓冲器中，并配置发送数据的位长。
        \item 置位 \hyperref[fielddesc:SPIUPDATE]{SPI\_UPDATE}，更新配置的寄存器值。
        \item 复位 PARLIO RX 模块。
        \item 配置 PARLIO RX 模块时钟。
        \item 关闭 PARLIO RX 模块 Core 时钟域。
        \item 将 PARLIO 接收模式配置为电平使能模式的子模式 1，配置 RX 模块数据总线宽度为 1 位，然后根据 SPI 的发送数据长度配置 \hyperref[fielddesc:PARLIORXDATABYTELEN]{PARL\_IO\_RX\_DATA\_BYTELEN}，最后置位 \hyperref[fielddesc:PARLIORXREGUPDATE]{PARL\_IO\_RX\_REG\_UPDATE}。
        \item 配置 PARLIO GDMA 接收链表。
        \item 置位 \hyperref[fielddesc:PARLIORXSTART]{PARL\_IO\_RX\_START}。
        \item 打开 PARLIO RX 模块 Core 时钟域。
        \item 置位 \hyperref[fielddesc:SPIUSR]{SPI\_USR}，开始传输 SPI 的数据。
        \item 轮询 GDMA SUC EOF 中断信号。
        \item 清除 \hyperref[fielddesc:PARLIORXSTART]{PARL\_IO\_RX\_START}。
    \end{itemize} 
    \item 实现 PARLIO 发送、SPI 接收，请参考以下操作过程：
    \begin{itemize}
        \item 配置 SPI 时钟。
        \item 将 SPI 配置为从机。
        \item 配置信号管脚，将 FSPICLK 管脚连接至 PAD\_CLK\_TX，FSPICS0 管脚连接至 TXD[16]，FSPID 管脚连接至 TXD[0]。
        \item 置位 \hyperref[fielddesc:SPIRDBITORDER]{SPI\_RD\_BIT\_ORDER}，翻转比特数据顺序。
        \item 置位 \hyperref[fielddesc:SPIUPDATE]{SPI\_UPDATE}，更新配置的寄存器值。  
        \item 复位 PARLIO TX 模块。
        \item 置位 \hyperref[fielddesc:PARLIOTXEOFINTCLR]{PARL\_IO\_TX\_EOF\_INT\_CLR} 和 \hyperref[fielddesc:PARLIOTXEOFINTENA]{PARL\_IO\_TX\_EOF\_INT\_ENA}。
        \item 配置 PARLIO TX 模块时钟。
        \item 关闭 PARLIO TX 模块 Core 时钟域。
        \item 配置数据总线宽度为 1 位，写 1 至 \hyperref[fielddesc:PARLIOTXHWVALIDEN]{PARL\_IO\_TX\_HW\_VALID\_EN}，然后配置 \hyperref[fielddesc:PARLIOTXBYTELEN]{PARL\_IO\_TX\_BYTELEN}。
        \item 配置 PARLIO GDMA 发送链表。
        \item 轮询 \hyperref[fielddesc:PARLIOTXREADY]{PARL\_IO\_TX\_READY}。        
        \item 写 1 至 \hyperref[fielddesc:PARLIOTXSTART]{PARL\_IO\_TX\_START}。
        \item 打开 PARLIO TX 模块 Core 时钟域。
        \item 开始数据传输。
        \item 轮询 \hyperref[fielddesc:PARLIOTXEOFINTST]{PARL\_IO\_TX\_EOF\_INT\_ST}。
        \item 置位 \hyperref[fielddesc:PARLIOTXEOFINTCLR]{PARL\_IO\_TX\_EOF\_INT\_CLR}。
        \item 关闭 PARLIO TX 模块 Core 时钟域。
        \item 清除 \hyperref[fielddesc:PARLIOTXSTART]{PARL\_IO\_TX\_START}。  
    \end{itemize} 
\end{itemize}

\subsection{与 I2S 进行数据传输}

在本示例中，外部 I2S 作为主机发送数据，PARLIO RX 模块作为从机接收数据。PARLIO 支持 I2S TDM MSB 对齐标准模式和 TDM PCM 标准模式下的传输。当 I2S 传输协议为 TDM MSB 对齐标准模式时，需将 PARLIO 的接收模式配置为电平使能模式。当 I2S 传输协议为 TDM PCM 标准模式时，需将 PARLIO 的接收模式配置为脉冲使能模式的子模式 10。

本节以 TDM PCM 标准模式为例，具体操作过程如下：

\begin{enumerate}
    \item 配置 I2S 时钟。
    \item 配置信号管脚，将 I2SO\_BCK\_out 连接至 PAD\_CLK\_RX，I2SO\_WS\_out 连接至 RXD[16]，I2SO\_Data\_out 管脚连接至 RXD[0]。
    \item 将 I2S 配置为主机。
    \item 配置所需 I2S TX 数据和通道模式，并置位 \hyperref[fielddesc:I2STXUPDATE]{I2S\_TX\_UPDATE}。
    \item 复位 I2S TX 模块和 TX FIFO。
    \item 使能 \hyperref[int:I2STXDONEINT]{I2S\_TX\_DONE\_INT}。 
    \item 配置 I2S GDMA 发送链表。
    \item 置位 \hyperref[fielddesc:I2STXSTOPEN]{I2S\_TX\_STOP\_EN}。
    \item 复位 PARLIO RX 模块。
    \item 配置 PARLIO RX 模块时钟。
    \item 关闭 PARLIO RX 模块 Core 时钟域。
    \item 将 PARLIO 接收模式配置为脉冲使能模式的子模式 10，配置 RX 模块数据总线宽度为 1 位，然后根据 I2S 的发送数据长度配置 \hyperref[fielddesc:PARLIORXDATABYTELEN]{PARL\_IO\_RX\_DATA\_BYTELEN}，最后置位 \hyperref[fielddesc:PARLIORXREGUPDATE]{PARL\_IO\_RX\_REG\_UPDATE}。
    \item 配置 PARLIO GDMA 接收链表。
    \item 置位 \hyperref[fielddesc:PARLIORXSTART]{PARL\_IO\_RX\_START}。
    \item 打开 PARLIO RX 模块 Core 时钟域。
    \item 置位 \hyperref[fielddesc:I2STXSTART]{I2S\_TX\_START}，开始发送数据。
    \item 轮询 \hyperref[int:I2STXDONEINT]{I2S\_TX\_DONE\_INT}。 
    \item 轮询 GDMA SUC EOF 中断信号。
    \item 清除 \hyperref[fielddesc:I2STXSTART]{I2S\_TX\_START}。
    \item 清除 \hyperref[fielddesc:PARLIORXSTART]{PARL\_IO\_RX\_START}。
\end{enumerate}

\subsection{与 LCD 进行数据传输}\label{parlio-lcd-data-trans}

\vspace{1em}
\begin{tiplisting}
    \chipname{} 不支持 LCD 接口。下文中 LCD 模块寄存器字段的详细描述，请参考相应 ESP 系列芯片的文档。
\end{tiplisting}
\vspace{-1em}

在本示例中，PPARLIO TX 模块作为主机向外部 LCD 发送数据（使用 I8080/MOTO6800 格式），设备间传输 8 位并行数据。具体操作过程如下：

\begin{enumerate}
    \item 配置信号管脚。将 CLK\_TX\_out 连接至 LCD 的像素时钟 PCLK，TXD[7:0] 连接至 LCD 数据输入脚，TXD[8] 连接至 LCD 的 CS 脚，TXD[9] 连接至 LCD 的 CD 脚。
    \item 复位 PARLIO TX 模块。
    \item 依次置位 \hyperref[fielddesc:PARLIOTXFIFOREMPTYINTCLR]{PARL\_IO\_TX\_FIFO\_REMPTY\_INT\_CLR}、\hyperref[fielddesc:PARLIOTXEOFINTCLR]{PARL\_IO\_TX\_EOF\_INT\_CLR}、\hyperref[fielddesc:PARLIOTXFIFOREMPTYINTENA]{PARL\_IO\_TX\_FIFO\_REMPTY\_INT\_ENA} 和  \hyperref[fielddesc:PARLIOTXEOFINTENA]{PARL\_IO\_TX\_EOF\_INT\_ENA}。
    \item 配置 PARLIO TX 模块时钟。
    \item 关闭 PARLIO TX 模块 Core 时钟域。
    \item 配置数据总线宽度为 16 位，然后配置 \hyperref[fielddesc:PARLIOTXBYTELEN]{PARL\_IO\_TX\_BYTELEN}。 
    \item 配置 GDMA 发送链表。注意，在链表中发送的数据应符合 I8080/MOTO6800 格式。低 8 位是有效的并行数据，第 9、10 位分别是 CS 和 CD，最高有效位是常数 1， 其余位为任意值。
    \item 轮询 \hyperref[fielddesc:PARLIOTXREADY]{PARL\_IO\_TX\_READY}。    
    \item 写 1 至 \hyperref[fielddesc:PARLIOTXSTART]{PARL\_IO\_TX\_START}。
    \item 打开 PARLIO TX 模块 Core 时钟域。
    \item 开始数据传输。
    \item 轮询 \hyperref[fielddesc:PARLIOTXEOFINTST]{PARL\_IO\_TX\_EOF\_INT\_ST}。
    \item 置位 \hyperref[fielddesc:PARLIOTXEOFINTCLR]{PARL\_IO\_TX\_EOF\_INT\_CLR}。
    \item 关闭 PARLIO TX 模块 Core 时钟域。
    \item 清除 \hyperref[fielddesc:PARLIOTXSTART]{PARL\_IO\_TX\_START}。 
\end{enumerate}

\section{中断} \label{sec:parlio-interrupts}


\begin{itemize}
\item \label{int:TXFIFOREMPTYINT}TX\_FIFO\_REMPTY\_INT：TX 模块 FIFO 空时触发该中断，此时 TX 发送的数据可能出现错误。
\item \label{int:RXFIFOWFULLINT}RX\_FIFO\_WFULL\_INT：RX 模块 FIFO 满时触发该中断，此时 RX 接收的数据可能出现错误。
\item \label{int:TXEOFINT}TX\_EOF\_INT：TX 模块发送完成完整的一帧数据时触发该中断。
\end{itemize}


% >>> If your module has no registers, delete sample text
%       for Sections Register Summary and Registers below <<<

%%%%%%%%%%%%%%%%%%%%%%%%
%%% Register Summary %%%
%%%%%%%%%%%%%%%%%%%%%%%%

% >>> Update the sample text below and insert
%     the info on register summary into the subfile <<<

\clearpage
\hypertarget{parlio-reg-summ}{}
\section{寄存器列表}
\label{sec:parlio-reg-summ}

本小节的所有地址均为相对于并行 IO 控制器基地址的地址偏移量（相对地址），具体基地址请见章节 \ref{mod:sysmem} \textit{\nameref{mod:sysmem}} 中的表 \ref{tab:sysmem-base-address}。

请查看章节 \hyperref[glossary-access-types]{\textit{寄存器的访问类型}}，了解“访问”列缩写的含义。

\subfile{\modulefiles/41-PARLIO-reg-summ__CN}


%%%%%%%%%%%%%%%%%
%%% Registers %%%
%%%%%%%%%%%%%%%%%

% >>> Update the sample text below and insert
%      the info on registers into the subfile <<<

\clearpage
\section{寄存器}

本小节的所有地址均为相对于并行 IO 控制器基地址的地址偏移量（相对地址），具体基地址请见章节 \ref{mod:sysmem} \textit{\nameref{mod:sysmem}} 中的表 \ref{tab:sysmem-base-address}。

\subfile{\modulefiles/41-PARLIO-reg__CN}

\end{document}
